\documentclass[11pt,a4paper]{article}

\usepackage{amsmath}
\usepackage{mathtools}
\usepackage{fontspec}
\usepackage{unicode-math}
\setmathfont{KpMath-Regular.otf}
\setmathfont{KpMath-Bold.otf}[version=bold]
\usepackage{microtype}
\usepackage[dvipsnames]{xcolor}
\usepackage[margin=24mm, top=28mm, bottom=28mm]{geometry}

\usepackage{fancyhdr}
\pagestyle{fancy}
\fancyhf{}
\fancyhead[L]{\small\itshape The Voxel \texorpdfstring{$T$}{T}-Matrix}
\fancyhead[R]{\small\itshape 19 September 2026}
\fancyfoot[C]{\thepage}
\renewcommand{\headrulewidth}{0.4pt}

\setcounter{secnumdepth}{3}
\usepackage{booktabs}
\usepackage{array}
\usepackage[font=small, labelfont=bf, labelsep=period]{caption}
\usepackage{setspace}
\setstretch{1.12}
\setlength{\parskip}{0.4em}
\usepackage[colorlinks=true, linkcolor=blue!60!black, urlcolor=blue!60!black]{hyperref}

\newcommand{\bI}{\mathbf{I}}
\newcommand{\bP}{\mathbf{P}}
\newcommand{\bT}{\mathbf{T}}
\newcommand{\bG}{\mathbf{G}}
\newcommand{\bu}{\mathbf{u}}
\newcommand{\bs}{\mathbf{s}}
\newcommand{\br}{\mathbf{r}}
\newcommand{\bx}{\mathbf{x}}
\newcommand{\bR}{\mathbf{R}}
\newcommand{\bpsi}{\boldsymbol{\psi}}
\newcommand{\tp}{^{\mkern-1.5mu\mathsf{T}}}

\title{\bfseries The Voxel \texorpdfstring{$T$}{T}-Matrix\\[4pt]
\large Discretising an elastic scatterer into cubes: the matched propagator,\\
and the accuracy of the result against exact Mie}
\author{}
\date{19 September 2026}

\begin{document}
\maketitle

\begin{abstract}
\noindent
A scatterer of arbitrary shape can be represented as a lattice of small cubes,
each carrying an analytic single-site $T$-matrix, coupled through the
elastodynamic propagator and solved by Foldy--Lax.  This note sets out that
construction, identifies the condition the propagator must satisfy for the
scheme to be consistent, and measures the accuracy that results.

The condition is a pairing.  A single-site $T$-matrix closes on a particular
functional of the internal field --- the value at the cube centre, or the
average over the cube --- and the propagator it is used with must close on the
same one.  The collocation closure requires the propagator averaged over the
\emph{receiver cell only}: one power of the cell form factor, not the two of a
Galerkin pairing and not the none of a point evaluation.

That average is given in closed form.  It is regular at every non-zero lattice
separation, its far field is a scalar multiple of the point propagator per
mode, and near contact it needs a quadrature order that follows the separation.

Accuracy is measured against the exact elastic Mie solution for a sphere, which
is the one arbiter that is both exact and laterally varying.  Two controls are
needed for the comparison to measure the method rather than the test: the
observation distance must reach the far field, and the staircase volume of the
voxelised body must be divided out.  With those in place the matched pairing is
$4.7$--$8.0\times$ closer to Mie than a point propagator in the Rayleigh
regime.  At $ka = 0.5$ the staircase of the sharp sphere decides the error and
the two propagators cannot be told apart there; on a smoothly graded sphere,
which has no staircase, the matched pairing converges at second order to the
exact solution at $k_Sa = 0.5$ and $1$, while the point propagator stalls near
$0.9\%$ however fine the grid.

\emph{These numbers supersede two earlier sets: one measured against an arbiter
that carried three compensating sign errors, and one whose far field counted
the incident wave's linear variation twice. \S\ref{sec:results} records what
changed.}
\end{abstract}

\tableofcontents

%=====================================================================
\section{The construction}\label{sec:construction}

\subsection{Foldy--Lax on a lattice of cubes}

Partition the scatterer into $N$ cubes of side $d$ on a regular lattice, with
centres $\br_m$.  Each cube is represented by a local $T$-matrix $\bT_{\rm loc}$
acting on a nine-component state --- three displacements and six Voigt strains,
\begin{equation}
  \bpsi = \big(u_z,\,u_x,\,u_y,\;\varepsilon_{zz},\,\varepsilon_{xx},\,
               \varepsilon_{yy},\,2\varepsilon_{xy},\,2\varepsilon_{zy},\,
               2\varepsilon_{zx}\big)\tp,
\end{equation}
which is the smallest state closed under the coupling: the source a cube
radiates depends on both the displacement and the strain it experiences.

The exciting field at cube $m$ is the incident field plus the field radiated by
every other cube,
\begin{equation}\label{eq:foldy}
  \bpsi_m = \bpsi^{\rm inc}_m + \sum_{n\neq m} \bP(\br_m-\br_n)\,\bT_{\rm loc}\,\bpsi_n,
  \qquad\text{that is}\qquad
  \big(\bI - \bP\bT\big)\bpsi = \bpsi^{\rm inc},
\end{equation}
with $\bP(\bs)$ the $9\times9$ propagator between cubes at separation $\bs$.
The composite $T$-matrix of the whole body follows by summing the sources,
$\bT_{\rm comp} = \sum_m \bT_{\rm loc}\bpsi_m$.

The self term is excluded from the sum: a cube's action on itself is contained
in $\bT_{\rm loc}$, which is where the singular part of the Green's tensor and
its static (Eshelby) content live.

\subsection{The propagator}

$\bP$ is assembled from the elastodynamic Green's tensor and its derivatives,
in blocks
\begin{equation}\label{eq:Pblocks}
  \bP = \begin{pmatrix} \bG & \mathbf{C}\\[2pt] \mathbf{H} & \mathbf{S}\end{pmatrix},
\end{equation}
where $\bG$ maps force to displacement, $\mathbf{S}$ maps stress source to
strain, and $\mathbf{C},\mathbf{H}$ are the cross couplings obtained by Voigt
contraction of $\partial G$ and $\partial\partial G$.  Written through the
Kupradze representation, every block is built from the two scalar Helmholtz
kernels $g_P$ and $g_S$ and their Cartesian derivative tensors, so the $P$ and
$S$ content is separable --- which \S\ref{sec:average} requires.

%=====================================================================
\section{The pairing condition}\label{sec:pairing}

\subsection{Two closures}

Volume-averaging the Lippmann--Schwinger equation with a uniform internal
ansatz gives
\begin{equation}\label{eq:double}
  \langle\nabla\bu\rangle = \langle\nabla\bu^0\rangle
   + \left[\frac{1}{V}\iint \partial\partial G(\bx-\bx')\,d\bx\,d\bx'\right]\Delta c\;E,
\end{equation}
a \emph{double} average, while evaluating the same equation at the cube centre
gives
\begin{equation}\label{eq:single}
  \nabla\bu(0) = \nabla\bu^0(0)
   + \left[\int_V \partial\partial G(-\bx')\,d\bx'\right]\Delta c\;E,
\end{equation}
a \emph{single} average.  Neither is wrong: one closes on the volume average of
the internal field, the other on its centre value.  For an ellipsoid they
coincide, because Eshelby's uniformity theorem makes the internal field
uniform; for a cube they differ.

\subsection{The condition}

A $T$-matrix built on \eqref{eq:single} must be used with a propagator that
also closes on the centre value.  The propagator between two cubes therefore
carries \emph{one} power of the cell form factor --- the average over the
receiver cell, with the source treated as a point --- and not two.  Writing the
cell indicator's transform,
\begin{equation}\label{eq:formfactor}
  \text{single average}\;\;\longleftrightarrow\;\;
    \operatorname{sinc}(k_xh)\operatorname{sinc}(k_yh)\operatorname{sinc}(k_zh),
  \qquad h = d/2,
\end{equation}
against the squared factor that a Galerkin pairing produces.  A point
propagator carries neither and is not a consistent partner for either closure.

\paragraph{Per mode.}  The form factor must be applied to each mode separately.
The $P$ and $S$ poles carry different $k_z$ at the same lateral wavenumber, so
a single scalar factor applied to the assembled $9\times9$ is incorrect.

%=====================================================================
\section{The single cell average}\label{sec:average}

\subsection{Definition and regularity}

For cubes at separation $\bs$ on the lattice, the matched propagator is
\begin{equation}\label{eq:avgdef}
  \langle\bP\rangle(\bs) = \frac{1}{d^3}\int_{|\mathbf{u}|_\infty\le h}
     \bP(\bs-\mathbf{u})\,d\mathbf{u}, \qquad h = d/2 .
\end{equation}
With a point source and a cubic receiver cell, the closest approach of the cell
to the source is its near face, so
\begin{equation}\label{eq:regular}
  |\bs-\mathbf{u}| \;\ge\; d/2 \qquad\text{for every lattice separation }\bs\neq0,
\end{equation}
and the integrand in \eqref{eq:avgdef} is regular everywhere on the domain.
The double average has no such property: its integrand reaches $r=0$ whenever
the cells touch.

\subsection{The far field, in closed form}

The midpoint error of a cell average is
\begin{equation}\label{eq:tail}
  \langle g\rangle - g = \frac{d^2}{24}\nabla^2 g + O(d^4),
\end{equation}
and for a Helmholtz kernel $\nabla^2 g_c = -\kappa_c^2 g_c$ exactly away from
the origin.  Derivatives commute with $\nabla^2$, so the same holds for every
derivative tensor and hence for every block of \eqref{eq:Pblocks}.  Therefore,
away from contact,
\begin{equation}\label{eq:scalartail}
  \langle\bP\rangle = \Big(1 - \frac{\kappa_c^2 d^2}{24}\Big)\bP + O(d^4)
  \qquad\text{per mode } c\in\{P,S\},
\end{equation}
the two modes carrying their own $\kappa$ and therefore not combinable before
the average is taken.

\subsection{Quadrature near contact}

Inside the near shell \eqref{eq:scalartail} is not accurate enough and
\eqref{eq:avgdef} is integrated directly.  The integrand is regular but peaked
near the face of the cell closest to the source, and a Gauss rule converges
slowly against a nearby peak.  The difficulty is a property of the separation
alone,
\begin{equation}\label{eq:qparam}
  q(\bs) = \frac{\big|\max(|\bs|-h,\,0)\big|}{d},
\end{equation}
componentwise in the clamp, which takes the values $0.5$ for a face neighbour,
$1/\sqrt2$ for an edge neighbour and $\sqrt3/2$ for a corner one.  Convergence
of a product Gauss rule against a $24$-point reference:

\begin{center}
\begin{tabular}{@{}lrrrr@{}}
\toprule
neighbour & $g=4$ & $g=8$ & $g=12$ & $g=16$\\
\midrule
face $(q=0.50)$   & $1.8\times10^{-2}$ & $2.9\times10^{-5}$ & $3.6\times10^{-8}$ & $1.4\times10^{-10}$\\
edge $(q=0.71)$   & $6.5\times10^{-4}$ & $9.0\times10^{-8}$ & $1.6\times10^{-10}$ & ---\\
corner $(q=0.87)$ & $6.1\times10^{-5}$ & $2.3\times10^{-10}$ & --- & ---\\
\bottomrule
\end{tabular}
\captionof{table}{The order should follow $q$: $16$ points per axis for face
neighbours, $12$ for edge, $8$ for corner, and the base order elsewhere.  Only
six and twelve separations respectively are expensive, so the adaptive rule
costs almost nothing.}
\end{center}

\noindent
Two independent evaluations of \eqref{eq:avgdef} --- direct quadrature and the
closed form \eqref{eq:scalartail} --- agree to $O(d^4)$, halving $d$ reducing
their difference by sixteen ($5.3\times10^{-4}$, $3.4\times10^{-5}$,
$2.1\times10^{-6}$ at $d = 0.2,\,0.1,\,0.05$), and the average reduces to the
point propagator at second order as \eqref{eq:tail} requires.

\subsection{Cost}

$\bP$ depends only on the separation, so on a regular lattice there are
$O(n^3)$ distinct blocks rather than $O(N^2)$ pairs.  Building the kernel once
per distinct separation --- which is what an FFT-accelerated solver does
anyway --- makes the averaged propagator affordable: the few near-contact
separations needing high-order quadrature are evaluated once, whatever the cell
count.

%=====================================================================
\section{Measuring the accuracy}\label{sec:measure}

The exact elastic Mie solution for a homogeneous sphere is the natural arbiter:
it is exact, and unlike a layered reference it is laterally varying.  Two
things must be controlled for the comparison to measure the discretisation
rather than the test.

\subsection{The observation distance}

If the scattered field is compared through a far-field asymptotic on one side
and an exact expression on the other, the neglected $1/(kr)$ term must be
small compared with the quantity being measured.  It is not a weak condition.
Measured at $ka = 0.1$ with the sphere divided $4$ per edge:

\begin{center}
\begin{tabular}{@{}rrr@{}}
\toprule
$r/a$ & $k_S r$ & volume-corrected error\\
\midrule
$100$ & $10$ & $0.5043$\\
$500$ & $50$ & $0.0575$\\
$5\times10^{3}$ & $500$ & $0.0075$\\
$5\times10^{4}$ & $5000$ & $0.0068$\\
\bottomrule
\end{tabular}
\captionof{table}{The comparison is usable only from $k_Sr\sim500$; at
$k_Sr=50$ it overstates the error sixfold.  It is not fully converged even at
the top: the last decade still moves the answer by $10\%$, which sets the
precision of every small number in \S\ref{sec:results}.  The earlier version of
this table read $0.0088$ at both of the last two rows and so looked converged;
correcting the arbiter made the errors smaller without making the harness
better, and the residual $1/(k r)$ term is now a visible fraction of them.}
\end{center}

\subsection{The staircase}

A cubic lattice does not fill a sphere, and its coverage is not monotonic in
the number of cells per edge.  For $n_{\rm sub} = 3,4,6$ the voxel volume
exceeds the sphere's by $26\%$, $4.7\%$ and $17\%$, and the raw error tracks
that dilution to within a factor $1.4$.  Two consequences follow.

First, the raw error must be divided by the volume ratio before it says
anything about the propagator.  Second, a refinement ladder in $n_{\rm sub}$
measures the geometry, not the method: the error moves non-monotonically
because the dilution does.  The comparison that isolates the discretisation is
at \emph{fixed} geometry --- the same voxel set, the same staircase, one
propagator exchanged for the other --- and its stability across resolutions is
what makes it meaningful.

%=====================================================================
\section{Results}\label{sec:results}

\begin{quote}
\small\itshape
What changed, and why.  The numbers of the first version compared two objects
that were \emph{both} wrong.  The exact Mie solution carried a spurious $(-1)^n$
in its multipole sum, inverting the density/dipole channel; the Foldy--Lax far
field carried a sign error in its force term; and the contrast-extraction route
carried a third error that cancelled the first.  They were found by an
impedance-null test and fixed together on 20 September 2026.  On 29 September a
second far-field defect was fixed: it counted the incident wave's linear
variation twice, an error of order $(ka)^2$ that no refinement removes.  The
$ka = 0.5$ ``reversal'' reported after the first fix (the averaged propagator
$0.71$--$0.84\times$ as good as the point one) did not survive it, and the
$ka = 1.5$ numbers of that version are withdrawn.  On 31 October 2026 the single
site became the uniform-strain closure exact at every frequency, in place of a
cube $T$-matrix with far-field form-factor corrections; on this sharp sphere that
moves the errors by under a tenth at $ka = 0.1$ and by up to a half at $ka = 0.5$,
without changing either verdict.
What survives throughout is \S\ref{sec:pairing}--\S\ref{sec:average}, which
concern the propagator alone and never passed through the arbiter.
\end{quote}

Volume-corrected far-field error against exact Mie, for a sphere with
$\Delta\lambda = 2$, $\Delta\mu = 1$, $\Delta\rho = 0.1$ on a background
$\alpha=5$, $\beta=3$, $\rho=2.5$ (\path{scripts/gate_sphere_cell_average_vs_mie.py}):

\begin{center}
\begin{tabular}{@{}rrrrr@{}}
\toprule
$ka$ & $n_{\rm sub}$ & point propagator & matched average & ratio\\
\midrule
$0.1$ & $3$ & $0.0056$ & $\mathbf{0.0012}$ & $4.66$\\
      & $4$ & $0.0072$ & $\mathbf{0.0010}$ & $7.36$\\
      & $6$ & $0.0087$ & $\mathbf{0.0011}$ & $8.01$\\
\midrule
$0.5$ & $3$ & $\mathbf{0.0055}$ & $0.0092$ & $0.60$\\
      & $4$ & $0.0098$ & $\mathbf{0.0028}$ & $3.48$\\
      & $6$ & $0.0084$ & $\mathbf{0.0062}$ & $1.35$\\
\bottomrule
\end{tabular}
\captionof{table}{Bold is the better of the pair.  In the Rayleigh regime the
matched average is $4.7$ to $8.0$ times closer to Mie, by a factor that varies
by $50\%$ across resolutions.  At $ka = 0.5$ the ratios straddle one: on a
staircased sphere no verdict is available there.}
\end{center}

\subsection{The pairing on a body without a staircase}

The smoothly graded sphere of radius $a = 10$\,m, its contrast full in the core
$r < a/2$ and falling to zero across the shell as a $C^2$ smoothstep, has no
staircase: the contrast and its slope vanish at the surface, so the voxelised
body converges to the body.  Solved with each cell at the profile's centre value
and the exact closure as its single site, against its exact solution, the far
field's error is (\path{scripts/measure_pairing_graded_sphere.py}):

\begin{center}
\begin{tabular}{@{}rrrrr@{}}
\toprule
 & \multicolumn{2}{c}{$k_Sa = 0.5$} & \multicolumn{2}{c}{$k_Sa = 1$}\\
\cmidrule(lr){2-3}\cmidrule(lr){4-5}
$n_{\rm sub}$ & point & matched average & point & matched average\\
\midrule
$6$  & $7.5\times10^{-3}$ & $2.8\times10^{-3}$ & $7.6\times10^{-3}$ & $1.6\times10^{-3}$\\
$8$  & $8.1\times10^{-3}$ & $1.6\times10^{-3}$ & $8.5\times10^{-3}$ & $1.3\times10^{-3}$\\
$12$ & $8.8\times10^{-3}$ & $1.8\times10^{-4}$ & $9.4\times10^{-3}$ & $1.6\times10^{-4}$\\
$16$ & $9.1\times10^{-3}$ & $3.5\times10^{-5}$ & $9.7\times10^{-3}$ & $3.5\times10^{-5}$\\
$24$ & $9.2\times10^{-3}$ & $1.3\times10^{-5}$ & $9.9\times10^{-3}$ & $1.2\times10^{-5}$\\
\bottomrule
\end{tabular}
\captionof{table}{With the point propagator the scheme does not converge: its
error stays near $0.9\%$ and grows slightly as the grid is refined.  With the
average over the receiving cell it converges at second order to the exact
solution, $700$ to $800$ times closer at $24$ cells across.}
\end{center}

\noindent
A plausible reading, not derived here: the two propagators differ at separation
$s$ by $(d^2/24)\nabla^2 G$ to leading order, and summed over the neighbours of a
cell the $\delta$-function content of $\nabla^2 G$ contributes an amount that does
not shrink with $d$.  Whatever the mechanism, the measurement shows that the
pairing of \S\ref{sec:pairing} is not a Rayleigh-regime refinement but the
condition for the scheme to converge to the continuum at all.

\subsection{The three regimes}

\paragraph{Rayleigh, $ka\lesssim0.1$.}  On the sharp sphere the matched average
reaches $0.10$--$0.12\%$, at the shape-factor floor of representing a sphere by
cubes, against $0.56$--$0.87\%$ for the point propagator.

\paragraph{Transition, $ka\sim0.5$ to $1$.}  On the sharp sphere the staircase
decides the error and the two propagators cannot be told apart; on the graded
sphere the matched pairing converges and the point propagator does not
(above).  The question the first version left open, whether the collocation
closure loses validity as $ka_{\rm sub}$ grows, is answered: at these
frequencies it does not.

\paragraph{Resonance, $ka\gtrsim1$.}  On a sphere with a sharp surface the error
is set by the representation of the curved boundary at a wavelength comparable
with the radius, not by the propagator.  The $ka = 1.5$ numbers of the first
version, $22$--$28\%$, were measured with the far field that counted the
incident wave twice and are not repeated here.

%=====================================================================
\section{Summary}

\begin{enumerate}
\item Partition the body into cubes; give each the single-site $T$-matrix of
      the uniform-strain closure, its self term exact at every frequency.
\item Couple them with the propagator averaged over the \emph{receiver} cell,
      \eqref{eq:avgdef}, applied per mode --- matching the collocation closure
      of the single-site $T$-matrix.
\item Use the closed form \eqref{eq:scalartail} beyond a near shell and direct
      quadrature inside it, with the order chosen from \eqref{eq:qparam}.
\item Solve \eqref{eq:foldy}; build the kernel once per distinct separation.
\item When comparing against an exact reference, reach the far field and
      divide out the staircase volume; compare at fixed geometry.
\end{enumerate}

\noindent
On a sharp sphere the method is accurate to about a tenth of a per cent in the
Rayleigh regime and is limited by the staircase of the curved surface beyond
it, where a representation that follows the boundary is required instead.

\noindent
The pairing of step~2 is the condition for the scheme to converge at all: on a
body without a staircase the matched pairing converges at second order to the
exact solution at $k_Sa = 0.5$ and $1$, while the point propagator stalls near
$0.9\%$.

\end{document}
